% 来源：OEBPS/Text/01_16.xhtml（OPF spine 顺序；仅第一册）
\cleardoublepage
\BookEntry[理查德·道金斯《自私的基因》]{理查德·道金斯《自私的基因》}

\BookCover{images/书单03自私的基因.jpg}{[英] 理查德·道金斯著| 卢允中等译| 中信出版社2012年版}

\begin{BookSummaryBox}
{\centering \textbf{全球最有影响力的公共知识分子之一\\ 道金斯的经典著作}\par}
\end{BookSummaryBox}

\subsubsection*{这本书谈什么？}

我们从哪里来，又将到哪里去？生命有何意义，我们该如何认识自己？这本书正是从这些问题出发，展开了论述。任何生物，包括我们自己，都只是求生的机器。在书中，复制、变异和淘汰三种简单机制可以演变出诸多生命现象，这就是认知科学的魅力。
\par

道金斯用简明通俗的形式，第一次把根据自然选择的社会学说这一部分充满趣味地介绍给大家。他提出：“我们生来是自私的。人类窥见了社会关系中基本的对称性和逻辑性，在我们有了更充分的理解之后，我们的政治见解会重新获得活力，从而对心理学的科学研究提供理论上的支柱。”当然，也正是在这个过程中，我们会对我们所受的苦难的根源有更多、更深的理解。
\par

\subsubsection*{芒格为什么谈起这本书？}

在芒格的生活、学习和决策方法中，关于查理的商业模式，以及这种伟大的商业模式是如何来的，我们可以从他推荐的阅读书目中看出一些端倪。像《枪炮、病菌与钢铁》《自私的基因》《冰河世纪》和《达尔文的盲点》都有一个共同的主题：关注前面提到的“竞争性毁灭”问题，研究为什么有些事物能够适应环境，存活下来，甚至在经过很长的时间之后占据统治地位。当这个主题被应用于投资选择时，芒格偏爱的企业就出现了：有些是通过消灭竞争对手而达到繁荣的企业（就像《自私的基因》里面描绘的），有些是通过合作而兴旺的企业（《达尔文的盲点》）。
\par

我们再次看到，查理能够熟练地应用许多学科的知识：有多少投资者能够像查理经常做的那样，考虑到如此之多、如此之复杂的因素呢？
\par

\subsubsection*{《自私的基因》作者理查德·道金斯如是说}

“成功基因的一个突出特性就是其无情的自私性。”
\par

“人类的记忆很脆弱，如果我们想要跟随我们的记忆，穿越时间，回到数百万年前，我们需要更加持久稳固的族群记忆。”
\par

“由于文化的出现，我们的人生的终极任务，将不只是繁殖，还有文化的创造和传承。”
\par

“我们具备足够的力量去抗拒我们那些与生俱来的自私基因。只有我们，我们人类，能够反抗自私的复制基因的暴政。”
\par

“信仰就是尽管缺乏证据，但仍然相信，甚至可能因为缺乏证据，所以相信。”
\par
